\documentclass[aspectratio=169]{beamer}
\input{header}
\coursecode{JKSSB}

\title{Current OS, Processes, Linux \& Utilities}
\subtitle{Lesson 60 --- Processes, Linux and Windows Utilities}
\author{JKSSB Computer Awareness}
\date{\today}

\begin{document}
\frame{\titlepage}

\begin{frame}{Deadlock and Its Four Necessary Conditions}
  \begin{itemize}
    \item A \key{deadlock} is a state in which a set of processes each wait
      for a resource held by another, so none can proceed.
    \item Four conditions must hold together for a deadlock to occur.
  \end{itemize}
  \begin{center}
    \begin{tabular}{p{0.26\textwidth}p{0.60\textwidth}}
      \toprule
      \textbf{Condition} & \textbf{Meaning} \\
      \midrule
      Mutual exclusion & A resource is held by only one process at a time. \\
      Hold and wait & A process holds one resource while waiting for others. \\
      No preemption & A resource cannot be forcibly taken away. \\
      Circular wait & Each process waits for a resource held by the next. \\
      \bottomrule
    \end{tabular}
  \end{center}
\end{frame}

\begin{frame}{Deadlock Recovery by Resource Preemption}
  \begin{itemize}
    \item Once a deadlock is detected, the system must \key{recover}.
    \item One recovery method is \key{resource preemption}: take a resource
      away from one process and give it to another so the cycle breaks.
    \item Preemption of a resource is thus a \key{recovery technique}, not
      one of the four necessary conditions.
  \end{itemize}
  \begin{alertblock}{Trap}
    \key{Preemption} is \bad{not} one of the four necessary conditions for
    deadlock. The condition is \key{no preemption}; preemption itself is used
    to recover. (PYQ-10)
  \end{alertblock}
\end{frame}

\begin{frame}{Priority Inversion}
  \begin{itemize}
    \item \key{Priority inversion} happens when a \key{high-priority} task is
      blocked because a \key{low-priority} task holds a resource it needs.
    \item While the low-priority task holds the resource, a
      \key{medium-priority} task can preempt it and run.
    \item The high-priority task therefore waits behind a lower-priority one
      --- the priority order is \bad{inverted}.
  \end{itemize}
  \begin{block}{One line}
    A high-priority task waits for a low-priority task, which is itself pushed
    aside by a medium-priority task.
  \end{block}
\end{frame}

\begin{frame}{Priority Inheritance: The Fix}
  \begin{itemize}
    \item \key{Priority inheritance}: while the low-priority task holds the
      resource, it \key{temporarily inherits} the high priority of the task
      waiting for it.
    \item It then runs at the inherited priority, releases the resource
      quickly, and drops back to its own priority.
    \item \key{Round-robin} scheduling by itself does \bad{not} prevent
      priority inversion. (PYQ-11)
  \end{itemize}
  \begin{block}{Keep the pair straight}
    Inversion = the problem; inheritance = the cure.
  \end{block}
\end{frame}

\begin{frame}{Starvation and Aging}
  \begin{itemize}
    \item \key{Starvation} (indefinite blocking): a process is ready to run
      but is \key{repeatedly passed over}, so it may wait forever.
    \item It commonly affects \key{low-priority} processes when
      higher-priority work keeps arriving.
    \item \key{Aging} is the standard cure: gradually \key{raise the priority}
      of a waiting process so it eventually runs.
  \end{itemize}
  \begin{alertblock}{Starvation vs deadlock}
    A starved process can still run in principle; a deadlocked process
    \bad{cannot} proceed at all. Starvation is a scheduling problem, deadlock
    is a resource-cycle problem.
  \end{alertblock}
\end{frame}

\begin{frame}{Preemptive Priority Scheduling}
  \begin{itemize}
    \item In \key{priority scheduling}, the ready process with the
      \key{highest priority} is chosen to run.
    \item It is \key{preemptive} when a running process can be
      \key{interrupted} the moment a higher-priority process becomes ready.
    \item It is \key{non-preemptive} when the running process keeps the CPU
      until it finishes or blocks.
  \end{itemize}
  \begin{block}{Watch out for}
    Strict preemptive priority scheduling tends to starve low-priority
    processes unless \key{aging} is added.
  \end{block}
\end{frame}

\begin{frame}{RTOS: Real-Time Operating System}
  \begin{itemize}
    \item \key{RTOS} stands for \key{Real-Time Operating System}.
    \item It is an operating system designed to respond within a
      \key{bounded time}, so that \key{deadlines} are met.
    \item \key{Hard real-time}: a missed deadline is a failure (a serious or
      catastrophic result).
    \item \key{Soft real-time}: a missed deadline only degrades quality; the
      result is still useful.
  \end{itemize}
  \begin{alertblock}{Trap}
    An RTOS is defined by \key{meeting time deadlines}, \bad{not} by being
    fast in general. A very fast system that misses deadlines is not
    real-time.
  \end{alertblock}
\end{frame}

\begin{frame}{The Linux Kernel}
  \begin{itemize}
    \item The \key{kernel} is the \key{core of the operating system}.
    \item It manages \key{processes} (scheduling), \key{memory},
      \key{devices} (through drivers), and \key{file systems}.
    \item It runs in \key{kernel space}; ordinary programs run in
      \key{user space} and request services through system calls.
    \item The Linux kernel is \key{open-source} (its source can be viewed,
      changed, and shared).
  \end{itemize}
  \begin{block}{Contrast}
    The \key{shell} is the command interpreter a user talks to; the
    \key{kernel} is the core that actually controls the hardware.
  \end{block}
\end{frame}

\begin{frame}{\texttt{/root} vs \texttt{/home}}
  \begin{center}
    \begin{tabular}{p{0.16\textwidth}p{0.70\textwidth}}
      \toprule
      \textbf{Path} & \textbf{Meaning} \\
      \midrule
      \texttt{/root} & The \key{home directory of the superuser} (the root user). \\
      \texttt{/home} & The parent directory that normally holds the home directories of \key{ordinary users} (for example \texttt{/home/alice}). \\
      \texttt{/} & The \key{root of the whole file system} (a single slash), different again from \texttt{/root}. \\
      \bottomrule
    \end{tabular}
  \end{center}
  \vspace{0.25cm}
  \begin{alertblock}{Trap}
    \texttt{/root} is the \bad{superuser's home}, not the top of the file
    system. The top of the file system is \texttt{/} (single slash).
  \end{alertblock}
\end{frame}

\begin{frame}{Windows and Linux File Systems}
  \begin{itemize}
    \item A \key{file system} organises how files and folders are stored and
      retrieved on a disk.
    \item \key{NTFS} = New Technology File System --- the \key{default for
      modern Windows}.
    \item \key{FAT32} = an older, widely compatible file system used for
      small drives and memory cards.
    \item \key{ext4} = the \key{native and common default file system for
      Linux}.
  \end{itemize}
  \begin{block}{Remember}
    NTFS $\rightarrow$ Windows; ext4 $\rightarrow$ Linux; FAT32 $\rightarrow$
    broad compatibility.
  \end{block}
\end{frame}

\begin{frame}{Windows Store}
  \begin{itemize}
    \item The \key{Windows Store} (now the \key{Microsoft Store}) is the
      built-in \key{online marketplace} in Windows.
    \item It is used to \key{search, download, and install} apps for Windows,
      and to keep them updated.
    \item Apps installed this way come from a \key{trusted catalogue} rather
      than from any website.
  \end{itemize}
  \begin{alertblock}{Trap}
    The Store is a \key{source of applications}, \bad{not} a disk-cleanup or
    defragmentation tool. Do not confuse it with the maintenance utilities.
  \end{alertblock}
\end{frame}

\begin{frame}{WordPad vs Notepad}
  \begin{center}
    \begin{tabular}{p{0.14\textwidth}p{0.36\textwidth}p{0.36\textwidth}}
      \toprule
      \textbf{Point} & \textbf{Notepad} & \textbf{WordPad} \\
      \midrule
      Text & \key{Plain text}, no formatting & \key{Basic formatting} (bold, italic, fonts, colour) \\
      Typical save & \texttt{.txt} & \texttt{.rtf} (rich text) \\
      Extra & Smallest, fastest text editor & Can insert pictures and objects \\
      \bottomrule
    \end{tabular}
  \end{center}
  \vspace{0.25cm}
  \muted{Notepad = plain text; WordPad = light rich text --- neither is a
  full word processor like MS Word.}
\end{frame}

\begin{frame}{Disk Cleanup}
  \begin{itemize}
    \item \key{Disk Cleanup} is a built-in Windows utility that
      \key{frees disk space}.
    \item It removes \key{unnecessary files}: temporary files, thumbnails,
      downloaded program files, and the \key{Recycle Bin} contents.
    \item You choose which categories to delete; the files are not part of
      your live documents.
  \end{itemize}
  \begin{block}{Use}
    When a drive is running low on space, Disk Cleanup deletes safe,
    non-essential files to recover capacity.
  \end{block}
\end{frame}

\begin{frame}{Defragmenter and the SSD Caution}
  \begin{itemize}
    \item A \key{hard disk drive (HDD)} stores a file in
      \key{fragments} spread over the disk.
    \item The \key{Disk Defragmenter} rearranges those fragments so each file
      is \key{contiguous}, which speeds up reading.
    \item \key{Caution}: do \bad{not} defragment a \key{solid-state drive
      (SSD)}.
  \end{itemize}
  \begin{alertblock}{Why not an SSD?}
    An SSD has \key{no moving parts}, so defragmenting gives no speed benefit
    and only adds \key{needless writes}, which shorten its life. Windows uses
    \key{TRIM} for SSDs instead. (Lesson scope)
  \end{alertblock}
\end{frame}

\begin{frame}{Imaging Fax: Windows Fax and Scan}
  \begin{itemize}
    \item \key{Windows Fax and Scan} is the built-in imaging and fax utility
      in Windows.
    \item It can \key{scan} a document (with a scanner) and
      \key{send or receive faxes} through a fax modem or fax service.
    \item Scanned items can be saved, viewed, and faxed as \key{images}.
  \end{itemize}
  \begin{block}{Remember}
    Imaging Fax in Windows = the \key{Windows Fax and Scan} app.
  \end{block}
\end{frame}

\begin{frame}{\texttt{xcopy /s}}
  \begin{itemize}
    \item \key{xcopy} is a Windows \key{command-line} command that copies
      files and directory trees.
    \item The switch \key{\texttt{/s}} copies the chosen directory
      \key{including its subdirectories}, but skips empty ones.
    \item Adding \key{\texttt{/e}} also copies \key{empty} subdirectories.
  \end{itemize}
  \begin{alertblock}{Exact term}
    Remember the command as written: \key{\texttt{xcopy /s}}. The \texttt{/s}
    means ``with subdirectories''.
  \end{alertblock}
\end{frame}

\begin{frame}{Lesson 60: Recall}
  \begin{itemize}
    \item Deadlock needs mutual exclusion, hold and wait, no preemption, and
      circular wait together.
    \item \key{Resource preemption} is a deadlock \key{recovery} method ---
      not one of the four conditions.
    \item \key{Priority inversion} = a high-priority task waits behind a
      low-priority holder; \key{priority inheritance} cures it.
    \item Round-robin by itself does \key{not} prevent priority inversion.
    \item \key{Starvation} = indefinite waiting; \key{aging} prevents it.
    \item Preemptive priority scheduling interrupts a running process for a
      higher-priority one.
    \item \key{RTOS} = Real-Time Operating System; defined by meeting
      deadlines (hard vs soft).
    \item The Linux \key{kernel} is the core; \texttt{/root} is the
      superuser's home, \texttt{/home} holds users' homes.
    \item NTFS = Windows; ext4 = Linux; FAT32 = compatibility.
    \item Notepad = plain text; WordPad = basic rich text; Disk Cleanup frees
      space; do not defragment an SSD; \texttt{xcopy /s} copies
      subdirectories.
  \end{itemize}
\end{frame}
\end{document}
